% Materiał podstawowy przejęty z AM I-12 z zachowaniem etykiet.
\chapter{Iloczyny nieskończone}
\label{ma1:ch:infinite-products}

W definicji szeregu przechodziliśmy do granicy sum częściowych. Analogicznie można badać granice iloczynów częściowych. Istnieje jednak ważna różnica: dodanie skończenie wielu składników nie może zmienić zbieżności szeregu, natomiast jeden czynnik równy zeru sprawia, że wszystkie dalsze iloczyny częściowe są zerowe. Odróżnimy zatem zbieżność do niezerowej granicy od zanikania iloczynów częściowych. Najpierw przeprowadzimy dowody bez odwoływania się do logarytmu, a następnie wykorzystamy funkcje elementarne z tomu I.

\section{Definicja i własności elementarne}
\begin{definicja}
\label{ma1:def:infinite-product}
Niech $K=\R$ lub $K=\C$, a $(u_n)_{n\ge1}$ będzie ciągiem w $K$. Iloczyn nieskończony $\prod_{n=1}^\infty u_n$ rozpatrujemy wraz z ciągiem iloczynów częściowych
\[
 P_0=1,\qquad P_N=\prod_{n=1}^N u_n\quad(N\in\N_+).
\]
Jeżeli wszystkie $u_n\ne0$, to iloczyn nazywamy zbieżnym, gdy istnieje granica $P_N\to P\in K\setminus\{0\}$. Wówczas piszemy $\prod_{n=1}^\infty u_n=P$. Jeżeli wszystkie czynniki są niezerowe, lecz nie istnieje skończona niezerowa granica $(P_N)$, iloczyn nazywamy rozbieżnym, również gdy $P_N\to0$. Jednoznaczność wartości iloczynu zbieżnego wynika z jednoznaczności granicy w przestrzeni metrycznej (tom I).
\end{definicja}
Przyjmujemy w tym rozdziale konwencję zbieżności do niezerowej granicy. Gdy występuje czynnik zerowy, wszystkie iloczyny częściowe od jego indeksu są zerowe; tę sytuację rozpatrujemy osobno. Można również stosować konwencję dopuszczającą skończenie wiele czynników zerowych i wymagającą zbieżności niezerowego ogona. W każdym przypadku trzeba odróżnić zwykłą granicę ciągu $(P_N)$ od zbieżności iloczynu w powyższym sensie.

\begin{twierdzenie}
\label{ma1:thm:product-elementary}
Niech $K=\R$ lub $K=\C$, a $(u_n)_{n\ge1},(v_n)_{n\ge1}$ będą ciągami w $K\setminus\{0\}$.
\begin{enumerate}[label=(\roman*)]
\item Dla dowolnego $m\in\N_+$, przy konwencji iloczynu pustego równego $1$, iloczyn $\prod_{n\ge1}u_n$ jest zbieżny wtedy i tylko wtedy, gdy $\prod_{n\ge m}u_n$ jest zbieżny. W przypadku zbieżności
$\prod_{n\ge1}u_n=(\prod_{n=1}^{m-1}u_n)(\prod_{n\ge m}u_n)$.
\item Jeżeli $\prod u_n=U$ i $\prod v_n=V$ są zbieżne, to $\prod(u_nv_n)=UV$ oraz $\prod u_n^{-1}=U^{-1}$.
\item Zbieżność $\prod u_n$ implikuje $u_n\to1$.
\end{enumerate}
\end{twierdzenie}
\begin{proof}
Oznaczmy $P_0=1$ oraz $P_N=\prod_{n=1}^N u_n$ dla $N\ge1$. Dla (i) iloczyn częściowy od $m$ do $N$ wynosi $P_N/P_{m-1}$, gdzie $P_{m-1}\ne0$. Jeśli $P_N\to P\ne0$, iloczyny od $m$ dążą do $P/P_{m-1}\ne0$. Odwrotnie, jeśli iloczyny od $m$ dążą do $T\ne0$, to $P_N=P_{m-1}\prod_{n=m}^N u_n\to P_{m-1}T\ne0$. Otrzymujemy również wzór na wartość iloczynu. Dla (ii) oznaczmy $Q_N=\prod_{n=1}^N v_n$. Iloczyny częściowe iloczynów o czynnikach $u_nv_n$ oraz $u_n^{-1}$ są odpowiednio $P_NQ_N$ i $1/P_N$. Rachunek granic daje granice $UV\ne0$ i $1/U\ne0$. Dla (iii) mamy $u_n=P_n/P_{n-1}$. Oba ciągi w ilorazie dążą do tego samego niezerowego $P$, zatem $u_n\to P/P=1$. W szczególności dla czynników rzeczywistych od pewnego indeksu $u_n>0$, ponieważ od pewnego indeksu $|u_n-1|<1/2$. Znak niezerowej granicy rzeczywistej jest znakiem iloczynu skończonego początku obejmującego wszystkie ujemne czynniki: dalsze iloczyny częściowe mają ten sam znak, a niezerowa granica zachowuje ten znak.
\end{proof}
\begin{przyklad}
\label{ma1:ex:telescoping-products}
Dla $n\ge2$ rozważmy trzy ciągi czynników (definicję iloczynu stosujemy po przeindeksowaniu). Z bezpośredniego skracania iloczynów skończonych dla $N\ge2$ otrzymujemy
\[
\begin{aligned}
 \prod_{n=2}^N(1-1/n)&=1/N,\\
 \prod_{n=2}^N(1+1/n)&=(N+1)/2,\\
 \prod_{n=2}^N(1-1/n^2)&=
 \prod_{n=2}^N\frac{n-1}{n}\frac{n+1}{n}
 =\frac{N+1}{2N}.
\end{aligned}
\]
Istotnie, w pierwszym iloczynie skracamy sąsiednie czynniki $(n-1)/n$, w drugim $(n+1)/n$, a trzeci jest iloczynem dwóch poprzednich, gdyż $(1-1/n)(1+1/n)=1-1/n^2$. W każdym przypadku czynniki dążą do $1$. W pierwszym przypadku iloczyny częściowe dążą do zera, w drugim do $+\infty$, a tylko w trzecim iloczyn jest zbieżny w sensie definicji \ref{ma1:def:infinite-product}, do $1/2$. Warunek konieczny zbieżności nie jest wystarczający.
\end{przyklad}

\section{Kryterium Cauchy'ego dla iloczynów}
\begin{twierdzenie}
\label{ma1:thm:product-cauchy}
Niech $(u_n)_{n\ge1}$ będzie ciągiem niezerowych liczb rzeczywistych lub zespolonych. Iloczyn $\prod_{n\ge1}u_n$ jest zbieżny wtedy i tylko wtedy, gdy
\[
 \forall\varepsilon\in\R_{>0}\ \exists N\in\N_+\
 \forall q,p\in\N_+\quad
 \left(q\ge p\ge N\ \Rightarrow\
 \left|\prod_{n=p}^{q}u_n-1\right|<\varepsilon\right).
\]
\end{twierdzenie}
\begin{proof}
Niech $P_0=1$, $P_n=\prod_{j=1}^n u_j$ dla $n\ge1$. Jeżeli $P_n\to P\ne0$, to od pewnego indeksu $|P_n-P|<|P|/2$, więc $|P_n|\ge|P|/2$. Dla $\varepsilon>0$ wybieramy $M$ tak, aby dla wszystkich $n\ge M$ zachodziły zarówno ta nierówność, jak i $|P_n-P|<\varepsilon|P|/4$. Gdy $q\ge p\ge M+1$, mamy
\[
 \left|\prod_{n=p}^q u_n-1\right|
 =\frac{|P_q-P_{p-1}|}{|P_{p-1}|}
 \le\frac{|P_q-P|+|P_{p-1}-P|}{|P|/2}
 <\varepsilon.
\]

Odwrotnie, stosujemy warunek z tezy najpierw dla $\varepsilon=1/2$. Istnieje $N_0\ge1$ takie, że dla każdego $q\ge N_0$ mamy $|\prod_{n=N_0}^q u_n-1|<1/2$. Ponieważ $P_{N_0-1}\ne0$ i $\prod_{n=N_0}^q u_n=P_q/P_{N_0-1}$, z nierówności trójkąta oraz nierówności odwrotnej otrzymujemy $1/2<|P_q/P_{N_0-1}|<3/2$. Stąd
\[
 \frac12|P_{N_0-1}|<|P_q|<\frac32|P_{N_0-1}|.
\]
Ciąg iloczynów częściowych jest więc ograniczony, również po uwzględnieniu skończonego początku; możemy przyjąć $B=\max(\{1,\frac32|P_{N_0-1}|\}\cup\{|P_n|:0\le n<N_0\})$. Wówczas $|P_n|\le B$ dla wszystkich $n\ge0$, także gdy $N_0=1$.

Dla dowolnego $\varepsilon>0$ wybieramy $N$ z warunku tezy dla dokładności $\varepsilon/B$. Dla $q>p\ge N$ otrzymujemy
\[
 |P_q-P_p|=|P_p|\left|\prod_{n=p+1}^q u_n-1\right|
 <B\frac{\varepsilon}{B}=\varepsilon.
\]
Przypadek $p=q$ jest trywialny, a dla $q<p$ zamieniamy indeksy. Z zupełności $\R$ i $\C$ (twierdzenie \ref{ma2:thm:finite}) ciąg $(P_n)$ ma granicę $P$. Dolne oszacowanie $|P_q|>|P_{N_0-1}|/2$ dla $q\ge N_0$ i zbieżność modułów dają $|P|\ge |P_{N_0-1}|/2>0$. Zatem iloczyn jest zbieżny.
\end{proof}

\section{Oszacowania skończonych iloczynów}
\begin{lemat}
\label{ma1:lem:product-finite-estimate}
Niech $m\in\N_+$ i niech $z_1,\ldots,z_m\in K$, gdzie $K=\R$ lub $K=\C$. Oznaczmy $s=\sum_{j=1}^m|z_j|$. Wówczas
\[
 \left|\prod_{j=1}^m(1+z_j)-1\right|
 \le \prod_{j=1}^m(1+|z_j|)-1.
\]
Jeżeli $s<1$, to
\[
 \prod_{j=1}^m(1+|z_j|)\le\frac1{1-s},
 \qquad
 \left|\prod_{j=1}^m(1+z_j)-1\right|
 \le\frac{s}{1-s}.
\]
\end{lemat}
\begin{proof}
Niech $I=\{1,\ldots,m\}$. Rozdzielność daje
\[
 \prod_{j=1}^m(1+z_j)=\sum_{E\subseteq I}\prod_{j\in E}z_j.
\]
Istotnie, z każdego czynnika wybieramy dokładnie jeden składnik: $z_j$ dla $j\in E$ lub $1$ dla $j\notin E$. Każdy podzbiór występuje dokładnie raz. Składnik odpowiadający zbiorowi pustemu wynosi $1$. Po odjęciu go i zastosowaniu nierówności trójkąta do pozostałych składników dostajemy pierwsze oszacowanie, gdyż suma ich modułów to dokładnie $\prod(1+|z_j|)-1$.

Niech $t_j=|z_j|\ge0$ oraz $s_k=\sum_{j=1}^k t_j$ dla $0\le k\le m$, gdzie $s_0=0$. Jeśli $s_m<1$, to $0\le s_k<1$ dla każdego $k$. Pokażemy indukcyjnie, że $\prod_{j=1}^k(1+t_j)\le1/(1-s_k)$. Dla $k=0$ zachodzi równość. W kroku indukcyjnym korzystamy z
\[
 (1-s_k)(1+t_k)
 =(1-s_{k-1}-t_k)(1+t_k)
 =1-s_{k-1}-s_{k-1}t_k-t_k^2
 \le1-s_{k-1}.
\]
Ponieważ oba mianowniki są dodatnie, wynika stąd $(1+t_k)/(1-s_{k-1})\le1/(1-s_k)$. Z założenia indukcyjnego, mnożąc przez dodatnie $1+t_k$, otrzymujemy
$\prod_{j=1}^k(1+t_j)\le(1+t_k)/(1-s_{k-1})\le1/(1-s_k)$, co kończy krok indukcyjny. Ostatnie oszacowanie wynika z pierwszego oraz $1/(1-s)-1=s/(1-s)$.
\end{proof}

\begin{lemat}
\label{ma1:lem:positive-product-bound}
Niech $m\in\N_+$ i niech $t_1,\ldots,t_m$ będą nieujemnymi liczbami rzeczywistymi. Wówczas
$\prod_{j=1}^m(1+t_j)\ge1+\sum_{j=1}^m t_j$.
Jeżeli ponadto $t_j\le1$ dla każdego $j$, to
$\prod_{j=1}^m(1-t_j)\ge1-\sum_{j=1}^m t_j$.
\end{lemat}
\begin{proof}
Pierwsza nierówność wynika z rozwinięcia iloczynu: składniki stopnia co najmniej dwa są nieujemne. Dla drugiej przeprowadzamy indukcję. Gdy $m=1$, zachodzi równość. Z założenia indukcyjnego, po pomnożeniu przez nieujemne $1-t_{m+1}$, mamy
\[
 \prod_{j=1}^{m+1}(1-t_j)
 \ge (1-\sum_{j=1}^m t_j)(1-t_{m+1})
 =1-\sum_{j=1}^{m+1}t_j+t_{m+1}\sum_{j=1}^m t_j
 \ge1-\sum_{j=1}^{m+1}t_j.
\]
\end{proof}

\section{Zbieżność bezwzględna}
\begin{definicja}
\label{ma1:def:absolute-product}
Niech $(a_n)_{n\ge1}$ będzie ciągiem liczb rzeczywistych lub zespolonych takim, że $1+a_n\ne0$ dla każdego $n$. Iloczyn $\prod_{n\ge1}(1+a_n)$ nazywamy bezwzględnie zbieżnym, jeżeli szereg $\sum_{n\ge1}|a_n|$ jest zbieżny.
\end{definicja}
\begin{twierdzenie}
\label{ma1:thm:absolute-product}
Niech $(a_n)_{n\ge1}$ będzie ciągiem w $\R$ lub $\C$, takim że $1+a_n\ne0$ dla każdego $n\ge1$. Jeśli szereg $\sum_{n=1}^\infty|a_n|$ jest zbieżny, to $\prod(1+a_n)$ jest zbieżny do niezerowej granicy $P$. Ponadto, oznaczając $P_N=\prod_{n=1}^N(1+a_n)$ i $P_0=1$, dla każdego całkowitego $N\ge0$ takiego, że $s_N=\sum_{n>N}|a_n|<1$, otrzymujemy
\[
 \left|\frac{P}{P_N}-1\right|\le\frac{s_N}{1-s_N}.
\]
\end{twierdzenie}
\begin{proof}
Ustalmy $\varepsilon>0$ i wybierzmy rzeczywiste $\delta$ z $0<\delta<\min(1/2,\varepsilon/2)$. Ponieważ szereg modułów jest zbieżny, istnieje $N_0$ z $\sum_{n\ge N_0}|a_n|<\delta$. Dla $q\ge p\ge N_0$ mamy $s=\sum_{n=p}^q|a_n|<\delta$. Z lematu \ref{ma1:lem:product-finite-estimate} wynika
\[
 \left|\prod_{n=p}^q(1+a_n)-1\right|
 \le\frac{s}{1-s}\le2s<\varepsilon.
\]
Kryterium \ref{ma1:thm:product-cauchy} daje zbieżność iloczynu do niezerowego $P$.

Dla ustalonego $N$ z $s_N<1$ oraz $q>N$ mamy $\sum_{n=N+1}^q|a_n|\le s_N$, a funkcja $s\mapsto s/(1-s)$ jest rosnąca na $[0,1)$, gdyż dla $0\le s\le t<1$ nierówność $s/(1-s)\le t/(1-t)$ po pomnożeniu przez dodatnie mianowniki jest równoważna $s\le t$. Z lematu otrzymujemy $|P_q/P_N-1|\le s_N/(1-s_N)$. Ponieważ $P_N\ne0$ i $P_q\to P$, rachunek granic oraz zbieżność modułów dają $|P_q/P_N-1|\to|P/P_N-1|$. Nieostra nierówność zachowuje się w granicy, co daje tezę.
\end{proof}
\begin{przyklad}
Dla $a_n=2^{-n}$ mamy $s_N=2^{-N}$, więc dla $N\ge1$ iloczyn $\prod_{n=1}^\infty(1+2^{-n})$ spełnia
$|P/P_N-1|\le1/(2^N-1)$. Dla $N=10$ jest to $1/1023<1/1000$. Możemy więc oszacować błąd względny względem obliczonego iloczynu częściowego bez znajomości dokładnej wartości $P$.
\end{przyklad}

\section{Czynniki o ustalonym znaku odchylenia od jedynki}
\begin{twierdzenie}
\label{ma1:thm:positive-product-series}
Niech $(t_n)_{n\ge1}$ będzie ciągiem nieujemnych liczb rzeczywistych.
\begin{enumerate}[label=(\roman*)]
\item Iloczyn $\prod_{n\ge1}(1+t_n)$ jest zbieżny wtedy i tylko wtedy, gdy szereg $\sum_{n\ge1}t_n$ jest zbieżny. Jeśli szereg jest rozbieżny, iloczyny częściowe dążą do $+\infty$.
\item Jeżeli dodatkowo $t_n<1$ dla wszystkich $n$, to iloczyn $\prod_{n\ge1}(1-t_n)$ jest zbieżny wtedy i tylko wtedy, gdy szereg $\sum_{n\ge1}t_n$ jest zbieżny. Jeśli szereg jest rozbieżny, iloczyny częściowe dążą do zera.
\end{enumerate}
\end{twierdzenie}
\begin{proof}
Zbieżność szeregu $\sum t_n$ daje w obu przypadkach zbieżność bezwzględną iloczynu, więc jego zbieżność wynika z twierdzenia \ref{ma1:thm:absolute-product}.

Załóżmy rozbieżność szeregu $\sum t_n$. Ponieważ $t_n\ge0$, jego sumy częściowe dążą do $+\infty$. Z lematu \ref{ma1:lem:positive-product-bound} mamy
$\prod_{n=1}^N(1+t_n)\ge1+\sum_{n=1}^N t_n$.
Prawa strona dąży do $+\infty$, co daje tezę w (i).

Dla (ii) zauważmy, że $1/(1-t_n)\ge1+t_n$: po pomnożeniu przez dodatnie $1-t_n$ otrzymujemy równoważną nierówność $1\ge1-t_n^2$. Stąd
\[
 0<\prod_{n=1}^N(1-t_n)
 \le\frac1{\prod_{n=1}^N(1+t_n)}
 \le\frac1{1+\sum_{n=1}^N t_n}.
\]
Dla każdego $\varepsilon>0$ z rozbieżności sum częściowych do $+\infty$ wynika, że od pewnego indeksu $1+\sum_{n=1}^N t_n>1/\varepsilon$. Ostatnia wielkość jest więc mniejsza od $\varepsilon$, a z podwójnego oszacowania iloczyny częściowe dążą do zera. Iloczyn nie jest zatem zbieżny w sensie przyjętej definicji.
\end{proof}
\begin{uwaga}
Twierdzenie daje równoważność zbieżności bezwzględnej iloczynu $\prod(1+a_n)$ oraz zbieżności iloczynu nieujemnych czynników $\prod(1+|a_n|)$, pod warunkiem że wszystkie czynniki wyjściowego iloczynu są niezerowe. Nie utożsamiamy tej własności ze zbieżnością iloczynu $\prod |1+a_n|$; na przykład dla $a_n=-2$ dla wszystkich $n$ mamy $\prod |1+a_n|=1$, podczas gdy $\prod(1+a_n)$ ma naprzemiennie iloczyny częściowe $-1$ i $1$, a $\sum|a_n|$ jest rozbieżny.
\end{uwaga}

\section{Zmiana kolejności czynników}
\begin{twierdzenie}
\label{ma1:thm:product-rearrangement}
Niech $(a_n)$ będzie ciągiem liczb rzeczywistych lub zespolonych, $1+a_n\ne0$ dla każdego $n$, i niech $\sum|a_n|$ będzie zbieżny. Dla każdej bijekcji $\pi:\N_+\to\N_+$ iloczyn $\prod_{k\ge1}(1+a_{\pi(k)})$ jest bezwzględnie zbieżny i ma tę samą wartość co $\prod_{n\ge1}(1+a_n)$.
\end{twierdzenie}
\begin{proof}
Z twierdzenia \ref{ma2:thm:absolute-rearrangement} szereg $\sum_{k\ge1}|a_{\pi(k)}|$ jest zbieżny. Z twierdzenia \ref{ma1:thm:absolute-product} oba iloczyny mają więc niezerowe granice; oznaczmy granicę iloczynu przestawionego przez $Q$, a granicę iloczynu wyjściowego przez $P$. Niech $P_N=\prod_{n=1}^N(1+a_n)$.
Niech $s_N=\sum_{n>N}|a_n|$. Wybierzmy $N$ tak duże, aby $s_N<1$. Z surjektywności $\pi$ istnieje $K_0=\max\{\pi^{-1}(1),\ldots,\pi^{-1}(N)\}$. Dla $K\ge K_0$ zbiór $F_K=\{\pi(1),\ldots,\pi(K)\}$ zawiera wszystkie indeksy od $1$ do $N$. Z przemienności skończonego mnożenia i iniektywności $\pi$ wynika
\[
 \prod_{k=1}^K(1+a_{\pi(k)})
 =P_N\prod_{n\in F_K\setminus\{1,\ldots,N\}}(1+a_n).
\]
Suma modułów $a_n$ po drugim zbiorze jest nie większa od $s_N$. Lemat \ref{ma1:lem:product-finite-estimate} daje zatem
\[
 \left|\frac{\prod_{k=1}^K(1+a_{\pi(k)})}{P_N}-1\right|
 \le\frac{s_N}{1-s_N}.
\]
Dla pustego zbioru drugi iloczyn wynosi $1$, więc oszacowanie również zachodzi. Przechodząc do granicy po $K$, otrzymujemy $|Q/P_N-1|\le s_N/(1-s_N)$. Teraz $P_N\to P\ne0$ i $s_N\to0$, więc $|Q/P-1|=0$, czyli $Q=P$.
\end{proof}

\section{Zbieżność warunkowa i związki z szeregami}
\begin{definicja}
\label{ma1:def:conditional-product}
Niech $(u_n)_{n\ge1}$ będzie ciągiem niezerowych liczb rzeczywistych lub zespolonych. Iloczyn $\prod_{n=1}^\infty u_n$ nazywamy zbieżnym warunkowo, jeżeli jest zbieżny do niezerowej granicy, lecz szereg $\sum_{n=1}^\infty|u_n-1|$ jest rozbieżny.
\end{definicja}
Nie można w ogólności zastąpić szeregu modułów w twierdzeniu \ref{ma1:thm:absolute-product} zwykłym szeregiem $\sum a_n$. Poniższe przykłady rozróżniają obie własności.

\begin{przyklad}
\label{ma1:ex:conditional-product}
Dla $k\ge1$ przyjmijmy $u_{2k-1}=1+1/k$ i $u_{2k}=k/(k+1)$. Wszystkie czynniki są dodatnie, a każdy iloczyn dwóch kolejnych czynników w parze wynosi $1$. Zatem
\[
 P_{2k}=1,\qquad P_{2k-1}=1+1/k.
\]
Oba podciągi dążą do $1$. Dla dowolnej dokładności $\varepsilon>0$ istnieje wspólny indeks $K$ taki, że oba rodzaje iloczynów częściowych są w odległości mniejszej niż $\varepsilon$ od $1$ dla $k\ge K$. Każdy $N\ge2K$ jest jednym z takich indeksów, więc $P_N\to1$. Jeśli $a_n=u_n-1$, to $a_{2k-1}=1/k$ oraz $a_{2k}=-1/(k+1)$. Szereg $\sum|a_n|$ jest rozbieżny, gdyż jego sumy częściowe o parzystych indeksach są nie mniejsze od $\sum_{k=1}^m1/k$ dla odpowiedniego $m$. Iloczyn jest zbieżny, lecz nie bezwzględnie, czyli jest zbieżny warunkowo. Zauważmy, że $\sum a_n$ jest w tym przykładzie również zbieżny: jego sumy częściowe o indeksie $2k$ wynoszą $\sum_{j=1}^k(1/j-1/(j+1))=1-1/(k+1)$, a sumy o indeksie $2k-1$ wynoszą $1-1/k+1/k=1$. Stąd suma szeregu wynosi $1$.
\end{przyklad}

\begin{przyklad}
\label{ma1:ex:conditional-product-rearranged}
Przestawmy czynniki z przykładu \ref{ma1:ex:conditional-product}, biorąc po dwa kolejne czynniki większe od $1$, a następnie jeden kolejny czynnik mniejszy od $1$. W $k$-tym bloku występują czynniki o indeksach $4k-3,4k-1,2k$. Każdy dodatni indeks parzysty ma jednoznacznie postać $2k$, a każdy dodatni indeks nieparzysty ma jednoznacznie jedną z postaci $4k-3$, $4k-1$, gdzie $k\ge1$. Te trzy rodziny są rozłączne, więc rzeczywiście otrzymujemy bijekcję indeksów.

Oznaczmy iloczyny częściowe przestawionego iloczynu przez $Q_N$. Po $m$ pełnych blokach, dla $m\ge1$, mamy
\[
 Q_{3m}=\prod_{j=1}^{2m}\frac{j+1}{j}
          \prod_{j=1}^{m}\frac{j}{j+1}
        =\frac{2m+1}{m+1}\longrightarrow2.
\]
Ponadto $Q_{3m+1}=Q_{3m}(1+1/(2m+1))$ i $Q_{3m+2}=Q_{3m+1}(1+1/(2m+2))$. Rachunek granic daje zbieżność obu tych podciągów do $2$. Dla dowolnej dokładności $\varepsilon>0$ wybieramy wspólny próg $M$ dla wszystkich trzech podciągów. Każde $N\ge3M$ ma postać $3m+r$, gdzie $m\ge M$ i $r\in\{0,1,2\}$, zatem $|Q_N-2|<\varepsilon$. Otrzymujemy iloczyn zbieżny do $2$, podczas gdy iloczyn wyjściowy był równy $1$. Założenia bezwzględnej zbieżności w twierdzeniu \ref{ma1:thm:product-rearrangement} nie można więc zastąpić samą zbieżnością.
\end{przyklad}

\begin{przyklad}
\label{ma1:ex:series-converges-product-zero}
Dla $k\ge1$ połóżmy $a_{2k-1}=1/\sqrt{k+1}$ i $a_{2k}=-1/\sqrt{k+1}$. Wszystkie czynniki $1+a_n$ są dodatnie. Parzyste sumy częściowe szeregu $\sum a_n$ wynoszą $0$, a nieparzyste $1/\sqrt{k+1}\to0$. Szereg jest zbieżny. Jednak
\[
 (1+a_{2k-1})(1+a_{2k})=1-\frac1{k+1}
 =\frac{k}{k+1},
 \qquad P_{2m}=\frac1{m+1}.
\]
Ponadto $P_{2m-1}=P_{2m-2}(1+1/\sqrt{m+1})=(1+1/\sqrt{m+1})/m\to0$. Wszystkie iloczyny częściowe dążą do zera, więc iloczyn jest rozbieżny w sensie definicji \ref{ma1:def:infinite-product}.
\end{przyklad}

\begin{przyklad}
\label{ma1:ex:product-converges-series-diverges}
Niech $t_k=1/\sqrt{k}$ dla $k\ge1$ oraz $u_{2k-1}=1+t_k$, $u_{2k}=1/(1+t_k)$. Ponownie $P_{2k}=1$, $P_{2k-1}=1+t_k\to1$, więc iloczyn jest zbieżny do $1$. Natomiast dla $a_n=u_n-1$ mamy
\[
 a_{2k-1}+a_{2k}=t_k-\frac{t_k}{1+t_k}
 =\frac{t_k^2}{1+t_k}\ge\frac1{2k},
\]
gdyż $0<t_k\le1$. Parzyste sumy częściowe szeregu $\sum a_n$ dążą do $+\infty$ z rozbieżności szeregu harmonicznego. Nieparzysta suma częściowa jest większa od poprzedniej parzystej o dodatni $t_k$. Stąd cały ciąg sum częściowych dąży do $+\infty$, mimo zbieżności iloczynu.
\end{przyklad}

\begin{uwaga}
Przykłady pokazują, że warunek $\sum a_n$ zbieżny nie jest ani konieczny, ani wystarczający do zbieżności $\prod(1+a_n)$. Poniżej wykorzystamy funkcje elementarne skonstruowane już w tomie I, aby uzyskać ścisły związek między iloczynami i szeregami logarytmów.
\end{uwaga}


\section{Logarytmy czynników rzeczywistych i zespolonych}
Korzystamy z funkcji wykładniczej, logarytmu rzeczywistego
oraz funkcji trygonometrycznych skonstruowanych w tomie I
(rozdział 18, Funkcje elementarne). W szczególności
$\exp(z+w)=\exp z\exp w$, funkcja $\exp$ jest ciągła
i nigdzie niezerowa, a $\log:(0,\infty)\to\R$ jest ciągłą
odwrotnością rzeczywistej funkcji wykładniczej.
Nie korzystamy z teorii różniczkowania funkcji zespolonych
ani z twierdzeń o szeregach potęgowych z dalszych rozdziałów.

\begin{twierdzenie}[Iloczyny dodatnie]
\label{ma2:thm:positive-log-product}
Dla ciągu dodatnich liczb rzeczywistych $(u_n)$
iloczyn $\prod_{n\ge1}u_n$ jest zbieżny wtedy i tylko wtedy,
gdy szereg $\sum_{n\ge1}\log u_n$ jest zbieżny.
W przypadku zbieżności
\[
 \prod_{n=1}^{\infty}u_n
 =\exp\left(\sum_{n=1}^{\infty}\log u_n\right).
\]
Jeżeli sumy częściowe szeregu logarytmów dążą do $+\infty$
lub $-\infty$, iloczyny częściowe dążą odpowiednio do
$+\infty$ lub $0$. Odwrotne implikacje również zachodzą.
\end{twierdzenie}
\begin{proof}
Dla każdego $N$ iloczyn częściowy $P_N$ jest dodatni,
a prawo logarytmowania iloczynu, zastosowane indukcyjnie,
daje $\log P_N=\sum_{n=1}^N\log u_n$.
Równoważnie $P_N=\exp(\sum_{n=1}^N\log u_n)$.
Jeżeli $P_N\to P\ne0$, to $P>0$ i z ciągłości logarytmu
sumy częściowe logarytmów dążą do $\log P$.
Jeśli sumy częściowe dążą do $S\in\R$, to z ciągłości
funkcji wykładniczej $P_N\to\exp S>0$.
Te same tożsamości oraz granice funkcji $\exp$ w obu
nieskończonościach i logarytmu w $0^+$ oraz $+\infty$
(tom I, rozdział 18) dowodzą pozostałych równoważności.
\end{proof}

Dla liczb zespolonych użyjemy wyłącznie logarytmu
określonego w prawej półpłaszczyźnie. Oznaczmy
\[
 L(w)=\log |w|+\mathrm i\arctan
 \frac{\operatorname{Im}w}{\operatorname{Re}w}
 \qquad(\operatorname{Re}w>0).
\]
Tutaj $\arctan$ przyjmuje wartości w $(-\pi/2,\pi/2)$.
Jest to funkcja ciągła, $L(1)=0$ i $\exp L(w)=w$.
Istotnie, jeśli $w=x+\mathrm i y$, $x>0$,
$\theta=\arctan(y/x)$, to $\cos\theta>0$ oraz
$\sin\theta=(y/x)\cos\theta$. Z tożsamości
$\sin^2\theta+\cos^2\theta=1$ dostajemy
$\cos\theta=x/|w|$ i $\sin\theta=y/|w|$.
Wzór Eulera z tomu I daje
$\exp L(w)=|w|(\cos\theta+\mathrm i\sin\theta)=w$.
Dla rzeczywistych $w>0$ mamy $L(w)=\log w$.

\begin{twierdzenie}[Kryterium logarytmiczne w przypadku zespolonym]
\label{ma2:thm:complex-log-product}
Niech $(u_n)$ będzie ciągiem niezerowych liczb zespolonych
takim, że $u_n\to1$. Wybierzmy $m\ge1$ tak, aby
$\operatorname{Re}u_n>0$ dla $n\ge m$.
Iloczyn $\prod_{n\ge1}u_n$ jest zbieżny wtedy i tylko wtedy,
gdy szereg $\sum_{n\ge m}L(u_n)$ jest zbieżny w $\C$.
W przypadku zbieżności
\[
 P=\left(\prod_{n=1}^{m-1}u_n\right)
       \exp\left(\sum_{n=m}^{\infty}L(u_n)\right).
\]
\end{twierdzenie}
\begin{proof}
Dla $N\ge m$ prawo dodawania argumentów funkcji
wykładniczej daje
\[
 \exp S_N=R_N,\qquad
 S_N=\sum_{n=m}^N L(u_n),\qquad R_N=\prod_{n=m}^N u_n .
\]
Zbieżność $S_N$ implikuje więc zbieżność $R_N$ do
niezerowej granicy. Twierdzenie
\ref{ma1:thm:product-elementary} pozwala dołączyć
skończony początek i daje wzór na $P$.

Odwrotnie, niech $R_N\to R\ne0$. Dla dostatecznie
dużego $N$ mamy $\operatorname{Re}(R_N/R)>0$,
a z ciągłości $L$ otrzymujemy $b_N=L(R_N/R)\to0$.
Wybierzmy $c\in\C$ z $\exp c=R$; istnienie takiej liczby
wynika z postaci trygonometrycznej niezerowej liczby
zespolonej (tom I, rozdział 18): jeśli
$R=|R|(\cos\theta+\mathrm i\sin\theta)$, przyjmujemy
$c=\log|R|+\mathrm i\theta$.
Wtedy $\exp(S_N-c-b_N)=1$.
Z prawa dodawania argumentów i wzoru Eulera wynika
$\exp z=1$ dokładnie wtedy, gdy $z\in2\pi\mathrm i\Z$:
dla $z=x+\mathrm i y$ moduł wykładniczej jest $\exp x$,
więc $x=0$, a $\cos y=1$ i $\sin y=0$ zachodzą
dokładnie dla całkowitych wielokrotności $2\pi$.
Istnieją zatem całkowite $k_N$ z
\[
 S_N=c+b_N+2\pi\mathrm i k_N .
\]
Dla kolejnych dostatecznie dużych indeksów mamy
\[
 2\pi\mathrm i(k_{N+1}-k_N)
 =L(u_{N+1})-(b_{N+1}-b_N)\longrightarrow0.
\]
Istotnie, $L(u_{N+1})\to L(1)=0$ oraz $b_N\to0$.
Lewa strona ma moduł równy zero lub co najmniej $2\pi$,
więc od pewnego indeksu $k_{N+1}=k_N$. Indukcyjnie
$k_N$ jest wtedy stały. Stąd $S_N$ ma skończoną granicę,
co kończy dowód.
\end{proof}

\section{Zbieżność bezwarunkowa iloczynów}
\begin{twierdzenie}[Kryterium iloczynów po zbiorach skończonych]
\label{ma2:thm:product-finite-tails}
Niech $(u_n)$ będzie ciągiem niezerowych skalarów.
Następujące warunki są równoważne:
\begin{enumerate}[label=(\roman*)]
\item każde przestawienie iloczynu jest zbieżne do
niezerowej granicy (nie zakładamy równości granic);
\item dla każdego $\varepsilon>0$ istnieje $N\ge1$ takie,
że dla każdego skończonego $H\subseteq\{N+1,N+2,\ldots\}$
zachodzi
\[
 \left|\prod_{n\in H}u_n-1\right|<\varepsilon;
\]
\item istnieje $P\ne0$ takie, że dla każdego $\varepsilon>0$
istnieje skończony zbiór $F_0\subseteq\N_+$, dla którego
$|\prod_{n\in F}u_n-P|<\varepsilon$ dla wszystkich
skończonych $F\supseteq F_0$.
\end{enumerate}
W przypadku spełnienia tych warunków wszystkie
przestawienia mają tę samą wartość $P$.
\end{twierdzenie}
\begin{proof}
Iloczyn po zbiorze pustym przyjmujemy równy $1$.
Gdyby (ii) nie zachodziło, istniałoby $\varepsilon>0$
takie, że poza każdym skończonym początkiem można wybrać
niepusty skończony zbiór $H$ z
$|\prod_{n\in H}u_n-1|\ge\varepsilon$.
Konstruujemy wyliczenie etapami jak w dowodzie twierdzenia
\ref{ma2:thm:unconditional-finite}: w każdym etapie
wybieramy taki zbiór poza maksimum dotychczas użytych
indeksów i numerem etapu, wypisujemy jego elementy
rosnąco, następnie wypisujemy najmniejszy dotychczas
niewypisany indeks. Indeksy nie powtarzają się, a po
$k$ etapach wypisano wszystkie indeksy $1,\ldots,k$
(indukcyjnie najmniejszy brakujący indeks jest wtedy
nie większy niż $k$). Otrzymujemy bijekcję.
Iloczyny po kolejnych wybranych blokach nie spełniają
warunku Cauchy'ego z twierdzenia
\ref{ma1:thm:product-cauchy}, ponieważ ich początkowe
pozycje w nowym wyliczeniu dążą do nieskończoności.
Iloczyn przestawiony jest rozbieżny, wbrew (i).
Zatem (i) implikuje (ii).

Załóżmy (ii). W szczególności dla kolejnych przedziałów
indeksów warunek Cauchy'ego daje zbieżność zwykłego
iloczynu do pewnego $P\ne0$. Oznaczmy jego iloczyny
częściowe przez $P_N$. Istnieje $B\ge1$ takie, że
$|P_N|\le B$ dla wszystkich $N$.
Dla $\varepsilon>0$ wybieramy $N$ tak duże, aby
$|P_N-P|<\varepsilon/2$ oraz warunek (ii) był spełniony
z dokładnością $\varepsilon/(2B)$.
Jeśli $F\supseteq\{1,\ldots,N\}$ jest skończony, to
\[
 \left|\prod_{n\in F}u_n-P\right|
 \le |P_N|\left|\prod_{n\in F\setminus\{1,\ldots,N\}}u_n-1\right|
      +|P_N-P|<\varepsilon.
\]
Otrzymujemy (iii).
Wreszcie dla dowolnej bijekcji $\pi$ każdy skończony
$F_0$ z (iii) jest zawarty w $\{\pi(1),\ldots,\pi(k)\}$
dla dostatecznie dużych $k$. Z (iii) iloczyny częściowe
przestawienia dążą do tego samego niezerowego $P$.
Dowodzi to (i) i równości wartości.
\end{proof}

\begin{uwaga}
Bezwzględna zbieżność iloczynu implikuje (ii):
jeżeli $s_N=\sum_{n>N}|u_n-1|<1$, to lemat
\ref{ma1:lem:product-finite-estimate} ogranicza
odchylenie od $1$ przez $s_N/(1-s_N)$, niezależnie
od wybranego skończonego zbioru indeksów.
Poniżej wykażemy również implikację odwrotną.
\end{uwaga}

\begin{lemat}
\label{ma2:lem:log-comparable}
Istnieje $\rho\in(0,1/2]$ takie, że dla $a\in\C$
z $|a|\le\rho$ zachodzi
\[
 \tfrac12|a|\le |L(1+a)|\le 2|a|.
\]
\end{lemat}
\begin{proof}
Dla $|a|\le1/2$ mamy $\operatorname{Re}(1+a)>0$.
Z ciągłości $L$ w $1$ i $L(1)=0$ wybieramy
$\rho\in(0,1/2]$ tak, aby $|L(1+a)|\le1/2$ dla
$|a|\le\rho$. Dla $b=L(1+a)$ zachodzi
$\exp b-1=a$. Z definicji wykładniczej
w tomie I, dla $|b|\le1/2$,
\[
 |\exp b-1-b|
 \le\sum_{k=2}^\infty\frac{|b|^k}{k!}
 \le |b|^2\sum_{k=2}^\infty (1/2)^{k-2}/2
 =|b|^2.
\]
Nierówność trójkąta i nierówność odwrotna dają
$\tfrac12|b|\le|a|\le\tfrac32|b|$, gdyż $|b|\le1/2$.
Stąd $\tfrac12|a|\le |b|\le2|a|$.
\end{proof}

\begin{twierdzenie}[Zbieżność bezwarunkowa a bezwzględna]
\label{ma2:thm:product-unconditional-absolute}
Dla niezerowych skalarów rzeczywistych lub zespolonych
$u_n$ iloczyn $\prod u_n$ jest bezwzględnie zbieżny wtedy
i tylko wtedy, gdy każde jego przestawienie jest zbieżne
do niezerowej granicy.
\end{twierdzenie}
\begin{proof}
Pierwsza implikacja wynika z twierdzenia
\ref{ma1:thm:product-rearrangement}.
Odwrotnie, załóżmy zbieżność wszystkich przestawień.
Zbieżność iloczynu w kolejności wyjściowej daje
$u_n\to1$. Wybierzmy $m$ tak duże, aby
$|u_n-1|\le\rho$ dla $n\ge m$, gdzie $\rho$ jest
z lematu \ref{ma2:lem:log-comparable}.
Każde przestawienie ogona jest zbieżne: poprzedzamy
je czynnikami $u_1,\ldots,u_{m-1}$, otrzymując
przestawienie całego iloczynu, po czym stosujemy
twierdzenie \ref{ma1:thm:product-elementary}.
Kryterium \ref{ma2:thm:complex-log-product} zastosowane
do każdego przestawienia ogona (po przeindeksowaniu)
daje zbieżność każdego przestawienia szeregu
$\sum_{n\ge m}L(u_n)$. W przypadku rzeczywistym możemy
traktować liczby rzeczywiste jako zespolone, bez zmiany
modułu i zbieżności.
Z twierdzenia \ref{ma2:thm:unconditional-finite} wynika
zbieżność bezwarunkowa szeregu logarytmów, a z twierdzenia
\ref{ma2:thm:finite-unconditional-absolute} jego
zbieżność bezwzględna, ponieważ $\C$ jest przestrzenią
skończenie wymiarową. Z lematu
\ref{ma2:lem:log-comparable} wynika
$|u_n-1|\le2|L(u_n)|$ dla $n\ge m$.
Kryterium porównawcze daje zbieżność
$\sum_{n\ge m}|u_n-1|$. Dołączenie skończonego początku
nie zmienia zbieżności szeregu norm.
\end{proof}

W następnym rozdziale przejdziemy do rodzin sumowalnych
i szeregów wielokrotnych.
